% ECONOMICS_PROBLEM: {"id": "E58-152", "title": "Liquidity backstops for nonbanks", "jel_code": "E58", "source_id": "OP-0326"}
% !TeX program = xelatex
\documentclass[11pt,a4paper]{article}
\usepackage[margin=23mm]{geometry}
\usepackage{fontspec}
\setmainfont{Latin Modern Roman}
\usepackage{amsmath,amssymb,mathtools}
\usepackage{xcolor,enumitem,booktabs,longtable,array,xurl,hyperref}
\definecolor{muted}{HTML}{53636C}
\hypersetup{colorlinks=true,urlcolor=blue,linkcolor=blue}
\setlength{\parindent}{0pt}
\setlength{\parskip}{5pt}
\newcommand{\R}{\mathbb R}
\newcommand{\E}{\mathbb E}
\newcommand{\N}{\mathbb N}
\newcommand{\ind}{\mathbf 1}
\newcommand{\Var}{\operatorname{Var}}
\newcommand{\Cov}{\operatorname{Cov}}
\newcommand{\supp}{\operatorname{supp}}
\DeclareMathOperator*{\argmin}{arg\,min}
\DeclareMathOperator*{\argmax}{arg\,max}
\newcommand{\entrymeta}[3]{{\sffamily\small\color{muted}\textbf{\texttt{#1}}\quad|\quad #2\par #3\par}}
\newcommand{\Problemref}[1]{\texttt{#1}}
\newcommand{\JELref}[2]{\texttt{#2} (JEL \texttt{#1})}
\begin{document}
\section*{Liquidity backstops for nonbanks}
\textbf{Problem ID:} \texttt{E58-152}\quad \textbf{JEL:} \texttt{E58}\par
% BEGIN ECONOMICS STATEMENT
\label{problem:OP-0326}
{\sffamily\small\color{muted}Original subject group: Banking, credit and financial stability\par}
\label{definitions:problem:30007608}
\textbf{Audit classification:} Explicit published open question. The BIS primary open-question section explicitly addresses central-bank intervention design; the Bank agenda supplies related scope. Collateral haircuts, fiscal accounting and endogenous buffers are expanded editorially.
\subsection*{Definitions and precise research target}
Fix nonbank institutions with assets, short-term funding, private liquidity buffers and potential access to central-bank lending. A facility \(a\) specifies eligible institutions, collateral, haircuts, interest charges, borrowing caps and activation conditions. Institution \(i\) chooses a pre-crisis buffer \(b_i(a)\) anticipating the facility. Let \(L(a,\omega;b(a))\) be crisis resource losses after this endogenous response, \(F(a,\omega)\) expected public losses, and \(C(a)\) routine operating costs. The task is to characterize
\[J(a)=E[L(a,\omega;b(a))+F(a,\omega)]+C(a),\qquad J_*=\inf_{a\in A}J(a).\]
Here \(A\) is the feasible facility set. Define shock state \(\omega\), its distribution, the institution perimeter and fiscal-risk accounting explicitly. Impose incentive and eligibility constraints ensuring institutions cannot profit from manufacturing collateral or misreporting solvency. Quantify the reduction in liquidity externalities and the offsetting loss of private self-insurance. Compare standing facilities with crisis-only lending while accounting for anticipated discretionary rescues. Determine when access must be coupled to supervision or prudential requirements. State which objects are identified from existing interventions and which require structural assumptions. A liquidity facility that prevents one historical fire sale does not establish an optimal access perimeter or demonstrate that moral hazard is negligible.

\textbf{Additional formal definitions and scope (editorial).} For collateral \(j\) of market value \(P_j q_j\), haircut \(h_j\in[0,1]\) permits at most \((1-h_j)P_j q_j\) of secured lending, subject to any separate borrowing cap. Buffer \(b_i(a)\) is an optimal private holding given anticipated access, with explicit private funding cost and losses in stress. Public fiscal cost is the unrecovered resource or distortionary-financing loss after collateral recovery; a loan disbursement alone is not a welfare loss. Specify repayment priority, maturity, central-bank balance-sheet limits and whether monetary-policy costs are included. The expectation in the objective uses a common shock law across facilities. If private choices or crisis prices are multiple, state a selector or minimize worst-case losses over the equilibrium correspondence. Standing-access and discretionary-rescue problems have different commitment assumptions. For a nonempty feasible set and a finite optimal value, the design target is an \(\varepsilon\)-optimal feasible rule for each specified \(\varepsilon>0\): its cost is at most the infimum plus \(\varepsilon\), or its welfare is at least the supremum minus \(\varepsilon\). Report infeasibility or an unbounded value separately. An exact optimizer is requested only under additional attainment assumptions; it need not exist for the general rule class.
\subsection*{Variants and assumption changes}
\begin{enumerate}
\item \textbf{Committed standing facility (source-motivated).} Announce access and haircuts before buffer choice and jointly choose prudential constraints; quantify private self-insurance changes.
\item \textbf{Discretionary stress lending (editorial).} Let authorities reoptimize after shocks with fiscal limits; distinguish the optimal crisis action from the ex ante effects of anticipating rescue.
\end{enumerate}
\subsection*{References and source audit}
\begin{enumerate}
\item BIS WP 972 cover verifies authors; October 2021 original, January 27, 2022 revision inspected; no DOI assigned. Locator: Section 5, printed pp.38--39 (PDF pp.41--42). Support: Nonbank regulation/backstop agenda. Full January 2022 revision accessible. URL: \url{https://www.bis.org/publications/working-paper-972-non-bank-financial-intermediaries-and-financial-stability.pdf}.
\item Official January 9, 2026 web version; locator: Central Bank Balance Sheet / Central banking toolkit, public-backstop paragraphs. Broad agenda; exact model editorial. URL: \url{https://www.bankofengland.co.uk/research/bank-of-england-agenda-for-research}.
\end{enumerate}
\begin{enumerate}\raggedright
\item\hypertarget{definitions:source:30007608:1}{} Sirio Aramonte; Andreas Schrimpf; Hyun Song Shin. 2021. \emph{Non-bank financial intermediaries and financial stability}. Section 5, printed pp.38--39 (PDF pp.41--42).
\url{https://www.bis.org/publications/working-paper-972-non-bank-financial-intermediaries-and-financial-stability.pdf}.
\item\hypertarget{definitions:source:30007608:2}{} Bank of England. 2026. \emph{Bank of England Agenda for Research, 2025–2028}. Central Bank Balance Sheet / Central banking toolkit, public-backstop paragraphs.
\url{https://www.bankofengland.co.uk/research/bank-of-england-agenda-for-research}.
\end{enumerate}

\subsection*{Common conventions: Source mathematical conventions}

These conventions apply to each entry unless explicitly replaced.

\textbf{Models and identification.} A primitive model \(m\) specifies all probability laws, feasible choices, information, equilibrium and measurement rules used by its target. The admissible class \(\mathcal M\) is fixed independently of outcomes. If \(P_O(m)\) is its observable probability law and \(\tau(m)\) the target, its identified set at observable law \(P\) is
\[
 \mathcal I_\tau(P)=\{\tau(m):m\in\mathcal M,\ P_O(m)=P\}.
\]
Point identification means that this set is a singleton. An empty set signals incompatibility of data and assumptions. Coordinate bounds need not describe the joint set. Bounds are sharp only if every retained target is attainable or arbitrarily approachable by compatible models. Finite bounds require explicit restrictions on unobserved quantities; unrestricted confounding, hidden wealth, extrapolation or arbitrary equilibrium selection can yield uninformative sets.

\textbf{Causal interventions.} A potential outcome \(Y_i(p)\) is the outcome under a complete policy regime \(p\), including timing, financing, information and market spillovers. The target population and units are fixed before treatment. With interference, use the complete assignment vector or a justified exposure mapping; individual no-interference notation alone cannot identify an economy-wide intervention. A difference of mean potential outcomes does not require the cross-world joint distribution. Conditioning on post-treatment migration, survival, enrollment or employment changes the estimand and needs separate justification.

\textbf{Statistical statements.} An estimator, confidence set, forecast or test specifies a sampling law, model class and loss/coverage criterion. Uniform \((1-\alpha)\)-coverage means \(\inf_{m\in\mathcal M}\Pr_m\{\tau(m)\in C_n\}\geq1-\alpha\), or an explicitly asymptotic version. The whole parameter space trivially has coverage, so informativeness requires a stated width, risk or power objective. A forecasting improvement is relative to a named benchmark, information set and evaluation population.

\textbf{Equilibrium and optimization.} An equilibrium comprises feasible plans, optimality, beliefs where needed, budget constraints and all market/resource clearing. The primitives or a stated benchmark define each correspondence \(\mathcal E(p;m)\). Nonemptiness is assumed or proved, never inferred just from compact policy sets. A selected-equilibrium target uses a declared selection rule; a robust target quantifies over all permitted equilibria. Use suprema or infima unless attainment is established. An \(\varepsilon\)-optimal policy is feasible and within \(\varepsilon\) of the finite optimum value. Report an empty feasible set separately.

\textbf{Welfare and accounting.} Normative utility, interpersonal normalization, social weights, numeraire, discounting and the use of public receipts are primitives. Behavioral utility need not coincide with the welfare criterion. Household welfare includes ownership income and rents once; adding profits again to compensating variation generally double-counts them. A monetary equivalent is defined by an explicit transfer and price convention, with monotonicity and range conditions. Average effects, mechanism contributions, transfers and welfare losses are different targets. Infinite sums require convergence and dynamic feasible plans require terminal or no-Ponzi conditions.

\textbf{Source versions.} A working-paper date, revision date and journal publication year are distinguished. A DOI for a journal version is not automatically the DOI of an earlier manuscript. Locators refer to the stated version and printed pages where specified. A metadata-only check does not verify a claim about a full-text conclusion. Every entry's source subsection narrows the provenance claim accordingly.
% END ECONOMICS STATEMENT
\end{document}
